\subsection{ Anéis}

\begin{enumerate}[label=10.\arabic*.]

    \item (A função $\varphi$ de Euler). Se $m$ e $n$ são inteiros $> 0$, e $m.d.c.(m,n) = 1$, prove que 
    $\varphi(mn) = \varphi(m)\varphi(n)$. Se $p$ é um primo, prove que $\varphi(p^n) = p^n - p^{n-1}$. \\
    (Sugestão: use o Teorema Chinês dos Restos para a primeira parte.) 

    \item ($\mathbb{Z}[\sqrt{2}]$ e $\mathbb{Q}(\sqrt{2})$). Verifique que os números da forma 
    $x + y\sqrt{2}, \quad x,y \in \mathbb{Z}$,formam um anel. Se $x,y \in \mathbb{Q}$, verifique 
    que eles formam um corpo. 

    \item (Sem divisores de zero). Demonstre que um anel associativo comutativo unitário finito 
    é um corpo se, e somente se, ele não possui divisores de zero. \\ 
    (Dica: Exercício 9.1.)

    \item (Sequências de Cauchy). Seja $F$ um corpo com um valor absoluto  $|\cdot| : F \to \mathbb{R}_{\geq 0}$
    (isto é, para quaisquer $x,y \in F$ valem: $|x|=0 \text{se, e somente se,} x=0$; $|xy| = |x||y|$; $|x+y| \leq |x|+|y|$). 
    Mostre que as sequências $(x_n)_{n \geq 0}$ sobre $F$ tais que, para cada $\varepsilon > 0$, existe um correspondente 
    $N > 0$ de modo que, para quaisquer $n,m \geq N$, vale $|x_n - x_m| < \varepsilon$, formam um anel com respeito às 
    operações de adição e multiplicação termo a termo. 

    \item ($R^(X)$). Sejam $X$ um conjunto e $R$ um anel. Verifique que o conjunto $R^{(X)}$ das funções $f: X \to R$ tais que $\{x \in X : f(x) \neq 0\}$ é finito é: 
    \begin{enumerate}[label=(\roman*)]
        \item um grupo com respeito à soma pontual de funções; 
        \item um anel definindo-se o produto por 
        \[
        (f*g)(x) = \sum_{y \in X} f(y)g(y^{-1}x), \quad \text{se $X$ for um grupo};
        \] 
        \item um espaço vetorial sobre $R$ sob a operação escalar $(rf)(x) = r f(x)$, se $R$ for um corpo.
    \end{enumerate}

    \item (Quaternions). Seja $\mathbb{H}$ o conjunto das matrizes complexas 
    $q = \begin{pmatrix} a+id & -b-ic \\ b-ic & a-id \end{pmatrix}, \quad a,b,c,d \in \mathbb{R}$.
    Prove que a soma e o produto matriciais usuais tornam $\mathbb{H}$ um anel de divisão (isto é, um 
    ``corpo possivelmente não comutativo''). 
    Prove que $q = a1 + b i + c j + d k$, com $1 = \begin{pmatrix} 1 & 0 \\ 0 & 1 \end{pmatrix}$,  
    $i = \begin{pmatrix} 0 & -1 \\ 1 & 0 \end{pmatrix}$, $j = \begin{pmatrix} 0 & -i \\ -i & 0 \end{pmatrix}$, e 
    $k = \begin{pmatrix} i & 0 \\ 0 & -i \end{pmatrix}$, e que $i^2 = -1 = j^2 e k = ij = -ji$.
    Prove que $q^2 = -1$ se $a = 0$ e $b^2 + c^2 + d^2 = 1$ 

\end{enumerate}
